\origpage{190}{ROBERTSON AND SEYMOUR}
\noindent of order, and the analogue of (4.3) holds. Indeed, this definition of the order of a separation extends to general matroids in the natural way, and again the analogue of (4.3) holds (with essentially the same proof). We suspect, but have not shown, that in a graph, Tutte-tangles and tangles are essentially the same objects. Some evidence for this lies in the fact that, for a connected planar graph, there is a 1-1 correspondence between its tangles and the tangles in a geometric dual [5].

\sectionhead{\textsc{References}}

1. G. A. Dirac, In abstrakten Graphen vorhandene vollst\"andige 4-Graphen und ihre Unterteilungen, \textit{Math. Nachr.} \textbf{22} (1960), 61--85.

2. N. Robertson and P. D. Seymour, Graph minors. IV. Tree-width and well-quasi-ordering, \textit{J. Combin. Theory Ser. B} \textbf{48} (1990), 227--254.

3. N. Robertson and P. D. Seymour, Graph minors. V. Excluding a planar graph, \textit{J. Combin. Theory Ser. B} \textbf{41} (1986), 92--114.

4. N. Robertson and P. D. Seymour, Graph minors. VIII. A Kuratowski theorem for general surfaces, \textit{J. Combin. Theory Ser. B} \textbf{48} (1990), 255--288.

5. N. Robertson and P. D. Seymour, Graph minors. XII. Circuits on a surface, submitted for publication.

6. N. Robertson and P. D. Seymour, Graph minors. XVII. Excluding a non-planar graph, submitted for publication.

7. N. Robertson, P. D. Seymour, and R. Thomas, Quickly excluding a planar graph, submitted for publication.

8. D. J. A. Welsh, ``Matroid Theory,'' Academic Press, London, 1976.
